\chapter{Fractions rationnelles}\label{ch:b1:fractions}

Une \omterm{def:b1:fractions:field}{fraction rationnelle} est un quotient de \omterm{def:b1:poly:def}{polynômes}. Le théorème
central de ce court chapitre --- la décomposition en éléments
simples\index{décomposition en éléments simples} --- casse un tel
quotient en une somme de briques élémentaires $\frac{c}{(X - a)^k}$.
Au-delà de son intérêt algébrique, c'est la machine standard pour
intégrer les fonctions rationnelles (\cref{ch:b1:integration}) et pour
sommer certaines séries (\cref{ch:b1:series}).

\section{Le corps $K(X)$}

\begin{definition}\label{def:b1:fractions:field}
Une \emph{fraction rationnelle}\index{fraction rationnelle} sur $K$ ($=
\R$ ou $\C$) est un quotient $F = \frac{A}{B}$ avec $A, B \in K[X]$, $B
\neq 0$~; deux quotients $\frac AB$ et $\frac{A'}{B'}$ sont identifiés
lorsque $AB' = A'B$. Toute fraction admet une \emph{forme irréductible}
avec $\gcd(A, B) = 1$, unique à une constante près. Muni des opérations
naturelles, l'\omterm{def:b1:logic:sets}{ensemble} $K(X)$ des fractions rationnelles est un \omterm{def:b1:structures:field}{corps}.

Les \emph{pôles}\index{pôle} de $F$ (sous forme irréductible) sont les
racines de $B$~; l'\emph{ordre} d'un pôle est sa \omterm{def:b1:poly:derivative}{multiplicité} comme
racine de $B$. Le \emph{degré} de $F$ est $\deg F = \deg A - \deg B \in
\Z \cup \{-\infty\}$.
\end{definition}

\begin{example}[Lire les pôles, les ordres et le degré]\label{ex:b1:fractions:vocab}
Soit $F = \dfrac{X^3 - X}{X^4 - 2X^3 + X^2}$. Factorisons les deux
étages~: numérateur $X(X-1)(X+1)$, dénominateur $X^2(X-1)^2$~; on
simplifie par le facteur commun $X(X-1)$~:
\[
F = \frac{X + 1}{X(X - 1)} \quad\text{(forme irréductible).}
\]
\omterm{def:b1:fractions:field}{Pôles}~: $0$ et $1$, tous deux \emph{simples} --- les ordres se lisent
sur le dénominateur irréductible, si bien que les racines doubles
apparentes du dénominateur initial n'ont aucune importance. Degré~:
$\deg F = 1 - 2 = -1$, visible asymptotiquement ($xF(x) \to 1$ quand $x
\to \infty$). Le degré se comporte comme le degré polynomial ($\deg FG
= \deg F + \deg G$, $\deg(F + G) \leq \max$), règle de comptabilité
utilisée constamment dans les chasses aux coefficients ci-dessous~:
chaque argument «~limite de $xF(x)$~» est un comptage de degrés
déguisé.
\end{example}

\begin{proposition}[Partie entière]\label{prop:b1:fractions:integerpart}
Tout $F = \frac AB \in K(X)$ s'écrit de manière unique $F = E + \frac
RB$ avec $E \in K[X]$ (la \emph{partie entière}, ou partie polynomiale,
de $F$) et $\deg R < \deg B$. On a $E \neq 0$ si et seulement si $\deg F
\geq 0$.
\end{proposition}

\begin{proof}
Division euclidienne $A = BE + R$ (\cref{thm:b1:poly:division}), divisée
par $B$. Unicité~: si $E + \frac RB = E' + \frac{R'}{B}$, alors $(E -
E')B = R' - R$ avec $\deg(R' - R) < \deg B$, ce qui force $E = E'$, puis
$R = R'$.
\end{proof}

\begin{example}[Simplifier d'abord, diviser ensuite]\label{ex:b1:fractions:integerpart}
Cherchons la partie entière de $F = \dfrac{X^3 + 1}{X^2 - 1}$. En
divisant à l'aveugle~: $X^3 + 1 = (X^2 - 1)X + (X + 1)$, donc $F = X +
\frac{X + 1}{X^2 - 1}$. Mais la fraction n'était pas irréductible~: $X^3
+ 1 = (X + 1)(X^2 - X + 1)$ et $X^2 - 1 = (X+1)(X-1)$ ont en commun le
facteur $X + 1$, et
\[
F = \frac{X^2 - X + 1}{X - 1} = X + \frac{1}{X - 1} :
\]
la même partie entière $X$, mais la partie fractionnaire se réduit à une
seule brique, et le «~\omterm{def:b1:fractions:field}{pôle}~» en $-1$ n'a jamais été un \omterm{def:b1:fractions:field}{pôle}. Toujours
réduire à la forme irréductible avant de chasser les \omterm{def:b1:fractions:field}{pôles}~: les \omterm{def:b1:fractions:field}{pôles}
de $F$ sont les racines du dénominateur \emph{irréductible}. (La partie
entière est insensible à la simplification, comme le garantit l'unicité
de la \cref{prop:b1:fractions:integerpart}.)
\end{example}

\section{Décomposition en éléments simples sur $\C$}

\begin{theorem}[Décomposition sur $\C$]\label{thm:b1:fractions:complex}
Soit $F = \frac AB \in \C(X)$ sous forme irréductible, avec $B = c\,(X -
a_1)^{m_1} \cdots (X - a_r)^{m_r}$. Alors $F$ s'écrit, de manière
unique,
\[
F = E + \sum_{i=1}^{r} \sum_{k=1}^{m_i}
\frac{c_{i,k}}{(X - a_i)^k},
\qquad E \in \C[X],\ c_{i,k} \in \C ,
\]
où $E$ est la partie entière de $F$.
\end{theorem}

\begin{proof}
D'après la \cref{prop:b1:fractions:integerpart}, on peut supposer $\deg
A < \deg B$ et démontrer la décomposition en somme avec $E = 0$.

\emph{Séparation des \omterm{def:b1:fractions:field}{pôles}.} Écrivons $B = (X - a_1)^{m_1} B_1$ avec
$B_1(a_1) \neq 0$. Les \omterm{def:b1:poly:def}{polynômes} $(X-a_1)^{m_1}$ et $B_1$ sont \omterm{cor:b1:arith:bezout}{premiers
entre eux} (aucune racine commune), donc, par Bézout dans $\C[X]$ (voir
la remarque du \cref{ch:b1:poly}), il existe $U_0, V_0$ tels que $U_0
(X-a_1)^{m_1} + V_0 B_1 = 1$~; en multipliant par $A$ et en posant $U =
AU_0$, $V = AV_0$, puis en divisant par $B$~:
\[
\frac AB = \frac{V}{(X - a_1)^{m_1}} + \frac{U}{B_1} .
\]
Les conditions de degré peuvent être imposées~: divisons $V$ par
$(X-a_1)^{m_1}$, disons $V = (X-a_1)^{m_1}Q + V_1$ avec $\deg V_1 <
m_1$, et absorbons le quotient dans le second terme ($U_1 = U +
QB_1$)~:
\[
\frac AB = \frac{V_1}{(X - a_1)^{m_1}} + \frac{U_1}{B_1},
\qquad \deg V_1 < m_1 ;
\]
la comparaison des degrés ($\deg A < \deg B$ et $\deg V_1 < m_1$) force
également $\deg U_1 < \deg B_1$. En itérant sur $\frac{U_1}{B_1}$, \omterm{def:b1:fractions:field}{pôle}
après \omterm{def:b1:fractions:field}{pôle}, on ramène tout au cas d'un seul \omterm{def:b1:fractions:field}{pôle} ci-dessous.

\emph{Un seul \omterm{def:b1:fractions:field}{pôle}.} Pour $\frac{A}{(X-a)^m}$ avec $\deg A < m$~:
développons $A$ suivant les puissances de $(X - a)$, $A =
\sum_{j=0}^{m-1} \alpha_j (X - a)^j$ (développement de Taylor d'un
\omterm{def:b1:poly:def}{polynôme}, comme dans la démonstration de la
\cref{prop:b1:poly:multiplicity})~; la division donne exactement les
briques $\frac{\alpha_j}{(X-a)^{m-j}}$. Concrètement, pour $\frac{X^2 +
1}{(X+2)^3}$~: en substituant $X = Y - 2$,
\[
X^2 + 1 = (Y - 2)^2 + 1 = Y^2 - 4Y + 5 ,
\qquad\text{donc}\qquad
\frac{X^2+1}{(X+2)^3} = \frac1{Y} - \frac4{Y^2} + \frac5{Y^3}
\]
avec $Y = X + 2$~: les trois briques apparaissent par simple division du
développement translaté --- la voie la plus rapide dès qu'un unique \omterm{def:b1:fractions:field}{pôle}
d'ordre élevé est en jeu, et celle que l'on recommande pour
l'\cref{exo:b1:fractions:3}.

\emph{Unicité.} Supposons que deux décompositions coïncident~; leur
différence est une identité $0 = \sum_{i,k} \frac{d_{i,k}}{(X -
a_i)^k}$. Multiplions par $(X - a_1)^{m_1}$~: chaque terme acquiert un
facteur s'annulant en $a_1$ \emph{sauf} celui d'indices $i = 1$, $k =
m_1$, dont le coefficient devient $d_{1,m_1}$ plus des termes portant au
moins un facteur $(X - a_1)$. L'évaluation en $a_1$ (légitime~: après la
multiplication, il ne reste plus de \omterm{def:b1:fractions:field}{pôle} en $a_1$) donne $d_{1,m_1} =
0$. Le coefficient de tête ayant disparu, on recommence avec $(X -
a_1)^{m_1 - 1}$, et ainsi de suite en descendant jusqu'à $k = 1$~; puis
on passe au \omterm{def:b1:fractions:field}{pôle} suivant. Tous les $d_{i,k}$ sont nuls~: la
décomposition est unique.
\end{proof}

\begin{method}[Calcul des coefficients]\label{met:b1:fractions:compute}
En pratique, on évite Bézout~; on combine~:
\begin{enumerate}
  \item la \emph{multiplication-évaluation} (ou «~méthode du cache~»)
        \emph{pour la puissance la plus
        haute}~: le coefficient de $\frac{1}{(X-a)^m}$ (où $m$ est
        l'ordre du \omterm{def:b1:fractions:field}{pôle} $a$) vaut
        \[
        c_{a,m} = \Bigl[\,(X - a)^m F\,\Bigr]_{X = a};
        \]
        pour un \omterm{def:b1:fractions:field}{pôle} \emph{simple} de $F = \frac AB$, c'est
        $\frac{A(a)}{B'(a)}$~;
  \item des \emph{évaluations} en des points commodes et des
        \emph{limites} de $xF(x)$ quand $x \to \infty$, qui fournissent
        des relations linéaires entre les coefficients restants~;
  \item les symétries de \emph{parité} ou de conjugaison, lorsqu'elles
        existent, pour diviser le travail par deux.
\end{enumerate}
\end{method}

\begin{remark}[Pièges courants avec les éléments simples]\label{rem:b1:fractions:pitfalls}
\begin{enumerate}
  \item \emph{Oublier la partie entière.} La décomposition en briques
        s'applique aux fractions de degré $< 0$~; lorsque $\deg A \geq
        \deg B$, il faut diviser d'abord
        (\cref{prop:b1:fractions:integerpart}), sinon la chasse aux
        coefficients produira des contradictions.
  \item \emph{Utiliser la multiplication-évaluation hors de son
        domaine.} Multiplier par $(X - a)^k$ puis évaluer en $a$ ne
        donne le coefficient que pour $k = m$, l'ordre \emph{plein} du
        \omterm{def:b1:fractions:field}{pôle}~; les coefficients d'ordre inférieur exigent d'autres
        relations (limites, évaluations) --- voir
        l'\cref{ex:b1:fractions:multiple}.
  \item \emph{Oublier de simplifier.} Les \omterm{def:b1:fractions:field}{pôles} se lisent sur la forme
        \emph{irréductible}~; un facteur commun au numérateur et au
        dénominateur crée des \omterm{def:b1:fractions:field}{pôles} fantômes
        (\cref{ex:b1:fractions:integerpart}).
  \item \emph{Mauvaise forme des briques sur $\R$.} Au-dessus d'un
        trinôme irréductible, les numérateurs sont \emph{affines}
        ($\alpha X + \beta$), non constants~; écrire $\frac{c}{X^2 + 1}$
        seul fait perdre des solutions --- les formes correctes sont
        dictées par le \cref{thm:b1:fractions:real}, jamais improvisées.
\end{enumerate}
\end{remark}

\begin{proof}[Démonstration de la formule du pôle simple]
Au voisinage d'un \omterm{def:b1:fractions:field}{pôle} simple $a$~: $B = (X - a) Q$ avec $Q(a) \neq 0$,
et $B' = Q + (X - a) Q'$, donc $B'(a) = Q(a)$. La valeur donnée par la
multiplication-évaluation est $\frac{A(a)}{Q(a)} =
\frac{A(a)}{B'(a)}$.
\end{proof}

\begin{example}\label{ex:b1:fractions:simple}
Décomposons $F = \dfrac{1}{X(X-1)(X-2)}$. Trois \omterm{def:b1:fractions:field}{pôles} simples~;
multiplication-évaluation en chacun~:
\[
c_0 = \frac{1}{(0-1)(0-2)} = \frac12,
\quad
c_1 = \frac{1}{1 \times (1 - 2)} = -1,
\quad
c_2 = \frac{1}{2 \times 1} = \frac12,
\]
donc $F = \dfrac{1/2}{X} - \dfrac{1}{X-1} + \dfrac{1/2}{X-2}$.
\emph{Vérification en $X = 3$~:} directement, $F(3) =
\frac{1}{3\cdot2\cdot1} = \frac16$~; par la décomposition,
$\frac{1/2}{3} - \frac{1}{2} + \frac{1/2}{1} = \frac16 - \frac12 +
\frac12 = \frac16$.
\end{example}

\begin{example}[Pôle multiple]\label{ex:b1:fractions:multiple}
Décomposons $F = \dfrac{X}{(X-1)^2 (X+1)}$. Forme~:
$\frac{a}{(X-1)^2} + \frac{b}{X-1} + \frac{c}{X+1}$.
\begin{itemize}
  \item Multiplication-évaluation au \omterm{def:b1:fractions:field}{pôle} double~: $a =
        \bigl[\frac{X}{X+1}\bigr]_{X=1} = \frac12$.
  \item Multiplication-évaluation en $-1$~: $c =
        \bigl[\frac{X}{(X-1)^2}\bigr]_{X=-1} = -\frac14$.
  \item Limite de $xF(x)$ en $\infty$~: $0 = b + c$, donc $b = \frac14$.
\end{itemize}
\[
F = \frac{1/2}{(X-1)^2} + \frac{1/4}{X-1} - \frac{1/4}{X+1} .
\]
\emph{Vérification en $X = 0$~:} $F(0) = 0$ et $\frac12 - \frac14 -
\frac14 = 0$.
\end{example}

\section{Décomposition sur $\R$}

\begin{theorem}[Décomposition sur $\R$]\label{thm:b1:fractions:real}
Soit $F \in \R(X)$ sous forme irréductible, de dénominateur
\[
B = c \prod_i (X - a_i)^{m_i} \prod_j (X^2 + p_j X + q_j)^{n_j}
\qquad (p_j^2 - 4q_j < 0).
\]
Alors $F$ se décompose de manière unique en sa partie entière, plus des
termes
\[
\frac{c_{i,k}}{(X - a_i)^k} \quad (1 \leq k \leq m_i),
\qquad
\frac{\alpha_{j,l}\, X + \beta_{j,l}}{(X^2 + p_j X + q_j)^{l}}
\quad (1 \leq l \leq n_j),
\]
à coefficients réels.
\end{theorem}

\begin{proof}
Décomposons sur $\C$ (\cref{thm:b1:fractions:complex}). Comme $F$ est
réelle, le coefficient au-dessus du \omterm{def:b1:fractions:field}{pôle} $\conj a$ (à chaque ordre) est
le \omterm{def:b1:complex:field}{conjugué} du coefficient au-dessus de $a$ (conjuguer la décomposition
et invoquer l'unicité). Regroupons chaque paire \omterm{def:b1:complex:field}{conjuguée}~:
\[
\frac{c}{(X - z)^l} + \frac{\conj c}{(X - \conj z)^l}
= \frac{c\,(X - \conj z)^l + \conj c\,(X - z)^l}{\bigl(X^2 - 2\Re(z)X
+ \abs z^2\bigr)^{l}},
\]
dont le numérateur est son propre \omterm{def:b1:complex:field}{conjugué}, donc réel, de degré $\leq
l$~; en lui retranchant des multiples du trinôme réel, on l'abaisse au
degré $\leq 1$ à chaque niveau $l$ (petite récurrence descendante). Les
\omterm{def:b1:fractions:field}{pôles} réels gardent leurs coefficients réels (la conjugaison les fixe).
L'unicité découle de l'unicité sur $\C$.
\end{proof}

\begin{example}[Le regroupement des conjugués en action]\label{ex:b1:fractions:pairing}
Le mécanisme de la démonstration, sur le plus petit cas~: sur $\C$, les
\omterm{def:b1:fractions:field}{pôles} de $\frac1{X^2+1}$ sont $\pm\iu$, de coefficients
$\frac1{2\iu}$ en $\iu$ et $\frac1{-2\iu}$ en $-\iu$ (par
multiplication-évaluation) --- \omterm{def:b1:complex:field}{conjugués} l'un de l'autre, comme le
théorème le prédit~:
\[
\frac{1}{X^2 + 1}
= \frac{1/(2\iu)}{X - \iu} - \frac{1/(2\iu)}{X + \iu} .
\]
En recombinant sur le dénominateur commun~:
\[
\frac{1}{2\iu}\cdot\frac{(X + \iu) - (X - \iu)}{X^2 + 1}
= \frac{1}{2\iu}\cdot\frac{2\iu}{X^2 + 1} = \frac{1}{X^2+1} :
\]
les parties imaginaires s'annulent et la brique réelle réapparaît
intacte. Pour des intégrandes réelles, on ne quitte en général jamais
$\R$ --- mais lorsqu'il s'agit d'évaluer des \emph{sommes} en des points
complexes (comme le fait le devoir maison avec les \omterm{def:b1:complex:unity}{racines de l'unité}),
les briques complexes sont la monnaie naturelle, et ce regroupement est
le taux de change entre les deux décompositions.
\end{example}

\begin{example}\label{ex:b1:fractions:real}
Décomposons $F = \dfrac{4}{(X^2+1)(X-1)^2}$ sur $\R$. Forme~:
$\frac{aX + b}{X^2 + 1} + \frac{c}{(X-1)^2} + \frac{d}{X - 1}$.
Multiplication-évaluation au \omterm{def:b1:fractions:field}{pôle} double~: $c =
\bigl[\frac{4}{X^2+1}\bigr]_{X=1} = 2$. Multiplication-évaluation au
\omterm{def:b1:fractions:field}{pôle} complexe $\iu$ (numérateur au-dessus de $X^2 + 1$ évalué via la
décomposition complexe, ou directement)~: on multiplie par $X^2 + 1$ et
on pose $X = \iu$~:
\[
a\iu + b = \frac{4}{(\iu - 1)^2} = \frac{4}{-2\iu} = 2\iu ,
\]
donc $a = 2$, $b = 0$. Limite de $xF(x)$ en l'infini~: $0 = a + d$, donc
$d = -2$. Ainsi
\[
F = \frac{2X}{X^2+1} + \frac{2}{(X-1)^2} - \frac{2}{X-1} .
\]
\emph{Vérification en $X = 0$~:} $F(0) = 4$ et $0 + 2 + 2 = 4$.
\end{example}

\begin{example}[Un dénominateur cubique, de bout en bout]\label{ex:b1:fractions:cubic}
Décomposons $F = \dfrac{1}{X^3 + 1}$ sur $\R$. Factorisons d'abord~:
$X^3 + 1 = (X + 1)(X^2 - X + 1)$, le trinôme ayant pour discriminant $-3
< 0$. Forme~: $\frac{a}{X+1} + \frac{bX + c}{X^2 - X + 1}$.
Multiplication-évaluation au \omterm{def:b1:fractions:field}{pôle} simple $-1$~: $a = \bigl[\frac1{X^2 -
X + 1}\bigr]_{X=-1} = \frac13$. Limite de $xF(x)$ en l'infini~: $0 = a +
b$, donc $b = -\frac13$. Évaluation en $X = 0$~: $1 = a + c$, donc $c =
\frac23$. Ainsi
\[
\frac{1}{X^3+1} = \frac13\Bigl(\frac{1}{X+1}
+ \frac{-X + 2}{X^2 - X + 1}\Bigr) ,
\]
ce que confirme $X = 1$~: membre de gauche $\frac12$, membre de droite
$\frac13\bigl(\frac12 + 1\bigr) = \frac12$. Notons l'économie~: trois
inconnues, trois faits linéaires peu coûteux (une
multiplication-évaluation, une limite, une évaluation), aucun
développement de quoi que ce soit --- le mode d'emploi de la
\cref{met:b1:fractions:compute} à l'état pur.
\end{example}

\begin{remark}[À quoi cela sert]
Une fois décomposée, une fonction rationnelle s'intègre terme à terme~:
les briques $\frac{1}{(x-a)^k}$ ont des primitives élémentaires, et les
briques $\frac{\alpha x + \beta}{(x^2 + px + q)^l}$ se ramènent à $\ln$
et $\arctan$ (\cref{ch:b1:integration}). Les sommes télescopiques sont
l'autre \omterm{def:b1:logic:map}{application} standard (\cref{exo:b1:fractions:8}).
\end{remark}

\begin{remark}[Où ce chapitre sert]\label{rem:b1:fractions:whereused}
Les éléments simples sont avant tout une étape de
\emph{prétraitement}~: le chapitre d'intégration
(\cref{ch:b1:integration}) fait passer chaque intégrande rationnelle par
le \cref{thm:b1:fractions:real} avant d'intégrer, et le chapitre sur les
séries (\cref{ch:b1:series}) télescope des termes rationnels exactement
comme dans l'\cref{exo:b1:fractions:8} et dans le devoir maison
ci-dessous --- lequel pousse la technique jusqu'à $\sum 1/k^2 =
\pi^2/6$. La dérivée logarithmique $P'/P = \sum m_i/(X - a_i)$
(\cref{exo:b1:poly:11}) réapparaît chaque fois que l'on étudie la
localisation des racines. Au-delà de ce volume, la décomposition de
$1/\chi(X)$ pour un \omterm{def:b1:diffeq:linear2}{polynôme caractéristique} $\chi$ sous-tend le calcul
des puissances de matrices et des transformées de Laplace dans le volume
de Licence 2~: les briques $\frac{c}{(X - a)^k}$ sont l'ombre algébrique
des solutions $t^{k-1}\eu^{at}$ rencontrées au \cref{ch:b1:diffeq}.
\end{remark}

\begin{omfigure}
\begin{tikzpicture}[xscale=1.75, yscale=0.42]
  \draw[->, thin] (-2.7,0) -- (2.8,0) node[below] {$x$};
  \draw[->, thin] (0,-4.6) -- (0,4.9) node[left] {$y$};
  \draw[dashed, gray] (-1,-4.4) -- (-1,4.6);
  \draw[dashed, gray] (1,-4.4) -- (1,4.6);
  \draw[omThm, thin] (-2.6,1.5) -- (2.7,1.5)
    node[above right, font=\small, pos=0.02] {$y = \lambda$};
  \draw[omDef, thick, domain=-2.6:-1.31, samples=80, smooth]
    plot (\x, {1/(\x+1) + 1/\x + 1/(\x-1)});
  \draw[omDef, thick, domain=-0.83:-0.20, samples=120, smooth]
    plot (\x, {1/(\x+1) + 1/\x + 1/(\x-1)});
  \draw[omDef, thick, domain=0.20:0.83, samples=120, smooth]
    plot (\x, {1/(\x+1) + 1/\x + 1/(\x-1)});
  \draw[omDef, thick, domain=1.31:2.7, samples=80, smooth]
    plot (\x, {1/(\x+1) + 1/\x + 1/(\x-1)});
  \fill[omThm] (-0.715,1.5) circle (0.045 and 0.19);
  \fill[omThm] (0.40,1.5) circle (0.045 and 0.19);
  \fill[omThm] (2.3,1.5) circle (0.045 and 0.19);
\end{tikzpicture}

{\small La fonction $F(x) = \frac1{x+1} + \frac1x +
\frac1{x-1}$, du type de l'\cref{exo:b1:fractions:12}~: strictement
décroissante sur chaque intervalle compris entre ses \omterm{def:b1:fractions:field}{pôles} $-1, 0, 1$
(en pointillés). Chaque niveau horizontal $\lambda > 0$ est franchi
exactement une fois par intervalle à droite du premier \omterm{def:b1:fractions:field}{pôle} (points
marqués)~: les solutions de $F = \lambda$ s'entrelacent avec les
\omterm{def:b1:fractions:field}{pôles}.}
\end{omfigure}

\section{Exercices}

\begin{exercise}[$\star$]\label{exo:b1:fractions:1}
Décomposer sur $\R$~: $\dfrac{1}{X^2 - 1}$~;
$\;\dfrac{X}{X^2 - 3X + 2}$~; $\;\dfrac{X^2 + 1}{X(X-1)}$
\emph{(attention à la partie entière)}.
\end{exercise}

\begin{exercise}[$\star$]\label{exo:b1:fractions:2}
Décomposer sur $\R$~: $\dfrac{1}{X(X^2 + 1)}$ et
$\dfrac{X^3}{X^2 + X + 1}$.
\end{exercise}

\begin{exercise}[$\star$]\label{exo:b1:fractions:3}
Décomposer $\dfrac{1}{X^2(X - 1)}$ et $\dfrac{X + 1}{(X - 1)^3}$
\emph{(pour la seconde, substituer $Y = X - 1$)}.
\end{exercise}

\begin{exercise}[$\star\star$]\label{exo:b1:fractions:4}
Décomposer sur $\C$, puis sur $\R$~: $\dfrac{1}{X^4 - 1}$.
\end{exercise}

\begin{exercise}[$\star\star$]\label{exo:b1:fractions:5}
Décomposer sur $\R$~: $\dfrac{X^2}{(X^2 + 1)^2}$, et en déduire une
primitive de $x \mapsto \dfrac{x^2}{(x^2+1)^2}$, sachant que
$\int \frac{\dd x}{(x^2+1)^2} = \frac12\bigl(\arctan x +
\frac{x}{x^2+1}\bigr) + C$.
\end{exercise}

\begin{exercise}[$\star\star$]\label{exo:b1:fractions:6}
Pour $n \in \N^*$, décomposer $F_n = \dfrac{n!}{X(X+1)\cdots(X+n)}$
\emph{(\omterm{def:b1:fractions:field}{pôles} simples en $0, -1, \dots, -n$~; utiliser la formule de
multiplication-évaluation et reconnaître des coefficients binomiaux)}.
\end{exercise}

\begin{exercise}[$\star\star$]\label{exo:b1:fractions:7}
En utilisant l'identité de l'\cref{exo:b1:poly:11} pour $P = X^n - 1$,
démontrer que
\[
\sum_{k=0}^{n-1} \frac{1}{X - \omega^k} = \frac{n X^{n-1}}{X^n - 1},
\qquad \omega = \eu^{2\iu\pi/n},
\]
et évaluer les deux membres en $X = 2$ pour $n = 4$ à titre de
vérification.
\end{exercise}

\begin{exercise}[$\star\star$]\label{exo:b1:fractions:8}
Décomposer $\dfrac{1}{k(k+1)(k+2)}$ et calculer
\[
S_n = \sum_{k=1}^{n} \frac{1}{k(k+1)(k+2)},
\qquad\text{puis}\qquad
\lim_{n \to \infty} S_n .
\]
\end{exercise}

\begin{exercise}[$\star\star\star$]\label{exo:b1:fractions:9}
Soit $P \in \R[X]$ unitaire de degré $n$, ayant $n$ racines réelles
distinctes $x_1 < \dots < x_n$. Démontrer que
\[
\sum_{i=1}^{n} \frac{1}{P'(x_i)} = 0 \quad (n \geq 2),
\qquad
\sum_{i=1}^{n} \frac{x_i^{\,n-1}}{P'(x_i)} = 1 .
\]
\emph{Indication~: décomposer $\frac{X^m}{P}$ pour $m \leq n - 1$ et
examiner le comportement des coefficients à l'infini --- ou utiliser
l'\omterm{thm:b1:poly:lagrange}{interpolation de Lagrange} (\cref{thm:b1:poly:lagrange}) du monôme
$X^m$ aux nœuds $x_i$.}
\end{exercise}

\begin{exercise}[$\star\star$]\label{exo:b1:fractions:10}
Décomposer $\dfrac{1}{X(X+1)^2}$ sur $\R$. En admettant la valeur
$\sum_{k \geq 1} \frac1{k^2} = \frac{\pi^2}6$ (démontrée dans le devoir
maison de ce chapitre), en déduire
\[
\sum_{n=1}^{\infty} \frac{1}{n(n+1)^2} = 2 - \frac{\pi^2}{6} .
\]
\end{exercise}

\begin{exercise}[$\star\star$]\label{exo:b1:fractions:11}
Décomposer sur $\R$~: $\dfrac{1}{(X^2+1)(X^2+4)}$, puis
$\dfrac{X^2}{(X^2+1)(X^2+4)}$. \emph{Indication~: les deux dénominateurs
sont des \omterm{def:b1:poly:def}{polynômes} en $X^2$~: décomposer d'abord
$\frac1{(Y+1)(Y+4)}$.}
\end{exercise}

\begin{exercise}[$\star\star\star$]\label{exo:b1:fractions:12}
(Équations séculaires) Soit $F = \sum_{i=1}^{r} \frac{c_i}{X - p_i}$
avec $p_1 < p_2 < \dots < p_r$ réels et tous les $c_i > 0$.
\begin{enumerate}
  \item Montrer que $F$ est strictement décroissante sur chaque
        intervalle de son domaine de définition, et donner ses limites
        en $\pm\infty$ et de part et d'autre de chaque \omterm{def:b1:fractions:field}{pôle}.
  \item En déduire que, pour tout $\lambda > 0$, l'équation $F(x) =
        \lambda$ a exactement $r$ solutions réelles, une dans chaque
        intervalle $\intoo{p_i}{p_{i+1}}$ et une au-delà de $p_r$.
        (De telles équations gouvernent les perturbations de valeurs
        propres~; le théorème des valeurs intermédiaires est utilisé
        ici au niveau du lycée et démontré au
        \cref{ch:b1:continuity}.)
\end{enumerate}
\end{exercise}

\section{Problème~: les éléments simples comme moteur}

\begin{problem}\label{pb:b1:fractions:1}
La décomposition en éléments simples a des airs de comptabilité~; ce
problème montre que c'est un moteur. Alimenté par la fraction
$\frac1{X(X+1)\cdots(X+k)}$, il télescope des familles entières de
sommes sous forme close~; alimenté par $\frac1{X^n - 1}$, il produit des
identités trigonométriques telles que
\[
\sum_{k=1}^{n-1}\frac{1}{\sin^2\frac{k\pi}{n}} = \frac{n^2-1}{3} ;
\]
et, poussée un cran plus loin, cette identité fait sortir l'une des
formules les plus célèbres des mathématiques, celle d'Euler~:
\[
\sum_{k=1}^{\infty}\frac1{k^2} = \frac{\pi^2}{6}
\]
--- obtenue ici sans rien d'autre que l'algèbre de ce chapitre et la
trigonométrie du lycée.\index{télescopage}\index{problème de Bâle}
Partout, $\omega = \eu^{2\iu\pi/n}$~; les limites de suites sont
utilisées au niveau du lycée (le \cref{ch:b1:seq} les formalise).

\medskip
\noindent\textbf{Partie I --- Le télescope.}
\begin{enumerate}
  \item Décomposer $\frac1{X(X+1)}$ et calculer exactement
        $\sum_{n=1}^{N} \frac1{n(n+1)}$~; en conclure que la somme tend
        vers $1$.
  \item Même chose pour $\frac1{X(X+2)}$~: montrer que $\sum_{n=1}^{N}
        \frac1{n(n+2)} = \frac12\bigl(\frac32 - \frac1{N+1} -
        \frac1{N+2}\bigr) \to \frac34$. (Avec un écart, \emph{deux}
        termes de bord survivent à chaque extrémité.)
  \item Formaliser le mécanisme~: si $F(X) = G(X) - G(X+1)$ pour une
        certaine \omterm{def:b1:fractions:field}{fraction rationnelle} $G$ sans \omterm{def:b1:fractions:field}{pôle} dans
        $\intco1{+\infty}$, alors $\sum_{n=1}^N F(n) = G(1) - G(N+1)$.
        Retrouver la valeur $\frac14$ de l'\cref{exo:b1:fractions:8} en
        exhibant le témoin $G$ pour $F = \frac1{X(X+1)(X+2)}$.
  \item Démontrer le télescope factoriel général~: pour $k \geq 1$,
        \[
        \frac{1}{X(X+1)\cdots(X+k)}
        = \frac1k\biggl(\frac{1}{X(X+1)\cdots(X+k-1)}
        - \frac{1}{(X+1)\cdots(X+k)}\biggr),
        \]
        et en déduire
        \[
        \sum_{n=1}^{\infty}\frac{1}{n(n+1)\cdots(n+k)}
        = \frac{1}{k \cdot k!} .
        \]
        Vérifier le cas $k = 2$ avec la question 3.
  \item Évaluer la décomposition de l'\cref{exo:b1:fractions:6} en des
        points bien choisis pour démontrer
        \[
        \sum_{j=0}^{n}(-1)^j\binom nj\,\frac1{j+1}
        = \frac1{n+1},
        \qquad
        \sum_{j=0}^{n}(-1)^j\binom nj\,\frac1{j+2}
        = \frac1{(n+1)(n+2)} .
        \]
\end{enumerate}

\medskip
\noindent\textbf{Partie II --- La fraction $1/(X^n - 1)$.}
\begin{enumerate}[resume]
  \item Montrer, par la formule de multiplication-évaluation, que
        \[
        \frac{1}{X^n - 1}
        = \frac1n\sum_{k=0}^{n-1}\frac{\omega^k}{X - \omega^k} .
        \]
  \item Deux vérifications~: vérifier la formule directement pour $n =
        2$, et montrer que la somme des $n$ coefficients est nulle pour
        $n \geq 2$ --- expliquer pourquoi elle doit l'être (considérer
        $xF(x)$ quand $x \to \infty$).
  \item Retrouver par multiplication-évaluation l'identité de
        l'\cref{exo:b1:fractions:7}~:
        $\frac{nX^{n-1}}{X^n-1} = \sum_k \frac1{X - \omega^k}$.
  \item Regrouper les \omterm{def:b1:fractions:field}{pôles} \omterm{def:b1:complex:field}{conjugués} pour démontrer la décomposition
        réelle~: avec $\theta_k = \frac{2k\pi}n$,
        \[
        \frac{\omega^k}{X - \omega^k} +
        \frac{\omega^{n-k}}{X - \omega^{n-k}}
        = \frac{2\cos\theta_k\,X - 2}
        {X^2 - 2\cos\theta_k\,X + 1} ,
        \]
        et écrire la décomposition réelle complète de
        $\frac1{X^n-1}$ (distinguer $n$ impair et $n$ pair).
  \item Spécialiser à $n = 4$ et confronter à
        l'\cref{exo:b1:fractions:4}.
\end{enumerate}

\medskip
\noindent\textbf{Partie III --- Sommes trigonométriques, et le
$\pi^2/6$ d'Euler.} Soit $P = 1 + X + \dots + X^{n-1}$, dont les racines
sont $\omega, \omega^2, \dots, \omega^{n-1}$ (toutes simples).
\begin{enumerate}[resume]
  \item En utilisant $\frac{P'}P = \sum_{k=1}^{n-1}\frac1{X - \omega^k}$
        (\cref{exo:b1:poly:11}), démontrer
        \[
        \sum_{k=1}^{n-1}\frac{1}{1 - \omega^k} = \frac{n-1}2 .
        \]
  \item Démontrer $\dfrac1{1 - \eu^{\iu\theta}} = \dfrac12 +
        \dfrac\iu2\,\frac{\cos(\theta/2)}{\sin(\theta/2)}$ pour
        $\theta \notin 2\pi\Z$ (factorisation par l'angle moitié,
        \cref{met:b1:complex:trig}), et déduire de la question 11
        que $\sum_{k=1}^{n-1}
        \frac{\cos(k\pi/n)}{\sin(k\pi/n)} = 0$ --- visible aussi
        par la symétrie $k \leftrightarrow n - k$.
  \item En dérivant l'identité de la question 11 (c'est-à-dire en
        utilisant $\bigl(\frac{P'}P\bigr)' = \frac{P''}P -
        \bigl(\frac{P'}P\bigr)^2$ évalué en $X = 1$), démontrer
        \[
        \sum_{k=1}^{n-1}\frac{1}{(1 - \omega^k)^2}
        = \frac{(n-1)(5-n)}{12} .
        \]
  \item En posant $\cot t = \frac{\cos t}{\sin t}$, déduire des
        questions 12--13 les deux formes closes
        \[
        \sum_{k=1}^{n-1}\cot^2\frac{k\pi}n = \frac{(n-1)(n-2)}3,
        \qquad
        \sum_{k=1}^{n-1}\frac1{\sin^2\frac{k\pi}n}
        = \frac{n^2-1}3 .
        \]
  \item Vérifier les deux formules à la main pour $n = 3$ et $n = 4$.
  \item Démontrer les inégalités $\cot t < \frac1t < \frac1{\sin t}$
        pour $t \in \intoo0{\frac\pi2}$ (à partir de $\sin t < t < \tan
        t$), et en déduire, pour $n = 2m + 1$ et $1 \leq k \leq m$~:
        \[
        \cot^2\frac{k\pi}n \;<\; \frac{n^2}{k^2\pi^2}
        \;<\; \frac1{\sin^2\frac{k\pi}n} .
        \]
  \item Sommer ces inégalités pour $k = 1, \dots, m$ (en utilisant la
        symétrie $k \leftrightarrow n - k$ pour diviser par deux les
        formules de la question 14) et encadrer~:
        \[
        \sum_{k=1}^{\infty}\frac1{k^2} = \frac{\pi^2}6 .
        \]
\end{enumerate}

\medskip
\noindent\textbf{Partie IV --- Briques d'ordre supérieur.}
\begin{enumerate}[resume]
  \item En élevant au carré la décomposition de $\frac1{X^2-1}$ et en
        redécomposant le terme croisé, démontrer
        \[
        \frac1{(X^2-1)^2}
        = \frac14\Bigl(\frac1{(X-1)^2} + \frac1{(X+1)^2}\Bigr)
        - \frac14\Bigl(\frac1{X-1} - \frac1{X+1}\Bigr),
        \]
        et le vérifier en $X = 0$.
  \item Combiner la question 18, le télescope et la valeur d'Euler
        (question 17) pour démontrer
        \[
        \sum_{n=2}^{\infty}\frac1{(n^2-1)^2}
        = \frac{\pi^2}{12} - \frac{11}{16} ,
        \]
        et confirmer la valeur numériquement à trois décimales.
  \item Dériver l'identité de la question 8 pour obtenir une forme
        close de $\sum_{k=0}^{n-1}\frac1{(X - \omega^k)^2}$, et la
        vérifier en $X = 2$, $n = 2$.
  \item Démontrer que, pour $k \geq 2$,
        \[
        \sum_{n=k}^{\infty}\frac{1}{\binom nk} = \frac{k}{k-1} .
        \]
        \emph{(Se ramener à la question 4 en écrivant $1/\binom nk$
        avec des factorielles.)}
  \item Pour $k = 3$, donner la somme partielle exacte
        $\sum_{n=3}^{N}\frac1{\binom n3}$ et sa limite.
\end{enumerate}

\medskip
\noindent\textbf{Partie V --- Synthèse.}
\begin{enumerate}[resume]
  \item En guise de couronnement, écrire la décomposition
        réelle complète de $\dfrac1{X^6 - 1}$ et la vérifier en $X =
        0$.
  \item Où exactement le problème a-t-il utilisé~: (i) l'unicité de la
        décomposition~; (ii) les \omterm{def:b1:complex:unity}{racines de l'unité} du
        \cref{ch:b1:complex}~; (iii) la dérivée logarithmique de
        l'\cref{exo:b1:poly:11}~? Une phrase pour chacune.
  \item Synthèse, en un court paragraphe~: une seule identité algébrique
        --- casser une fraction en briques --- a engendré des sommes
        exactes, des identités trigonométriques et $\pi^2/6$. Commenter
        la répartition des tâches entre l'algèbre (décompositions
        exactes, valables partout) et l'analyse (limites, encadrements),
        et indiquer où chaque fil est industrialisé~: télescopage et
        comparaison au \cref{ch:b1:series}, intégration des briques au
        \cref{ch:b1:integration}.
\end{enumerate}
\end{problem}
